\documentclass[11pt]{article}
\usepackage[T1]{fontenc}
\usepackage[margin=1in]{geometry}
\usepackage{amsmath,amssymb,amsthm,booktabs,array}
\usepackage[hidelinks]{hyperref}
\newtheorem{theorem}{Theorem}
\newtheorem{lemma}{Lemma}
\newtheorem{corollary}{Corollary}
\newcommand{\ZZ}{\mathbb Z}
\newcommand{\QQ}{\mathbb Q}
\newcommand{\FF}{\mathbb F}
\newcommand{\ev}{\operatorname{ev}}
\title{A jet-space proof and extension of Su's theorem on Bala's\ totient-square period conjectures}
\author{Haobo Ma\\\small ChronoAI Pte Ltd, Singapore
  \\[0.5em]Wenlin Zhang\\\small National University of Singapore
  \\\small ChronoAI Pte Ltd, Singapore}
\date{14 November 2026\\\small CC BY 4.0}
\begin{document}
\maketitle
\begin{abstract}
Using finite jets, we give an alternative proof of the totient-square period conjecture in
Peter Bala's comment of 16 February 2026 on OEIS A281181, revision 103,
previously proved by Lien-Hung Su (2026, Theorem 3.2 and Corollary 3.3),
and extend the argument to every positive integer flow exponent.
If $C=\cosh H$, $H'=C^3$ and $H(0)=0$, then the integers
$a(n)=(2n)![t^{2n}]C(t)$ satisfy
$a(n+\varphi(k)^2)\equiv a(n)\pmod{k}$ for every $k\geq1$ and $n\geq k^2$.
A characteristic-$p$ derivation identity on truncated polynomial rings,
followed by a commuting-matrix congruence, gives a self-contained proof.
The same argument covers every nonnegative power of the series $C$
for every positive integer exponent in $H'=C^q$, and hence four companion
OEIS conjectures and further even-exponent specializations.
Exact finite-dimensional certificates also recover
the source's small-modulus examples.
\end{abstract}

\section{The source and the statement}
Bala's conjecture asks for eventual periodicity modulo every positive
integer, with period dividing the square of Euler's totient
\cite{oeis181}. Here $\varphi(k)$ counts the units modulo $k$, with
$\varphi(1)=1$; saying that a period divides $T$ means that shifting by
$T$ preserves the sequence on a tail. Su already proves these conjectures
and the stronger shift $\varphi(k)^2/2$ for $k\geq3$, with sharper
preperiod bounds \cite[Theorem 3.2, Corollary 3.3]{su}.
The present proof uses a polynomial annihilator rather than a decomposition
into its invertible and nilpotent parts. Its bound $k^2$ is convenient,
not an improvement on Su's bounds. The inverse-function derivative
representation is also recorded by Bala \cite[Corollary 1]{bala}.
The extension here removes the oddness restriction on $q$ for the
square-totient shift. In particular it covers the even-exponent companion
flows, but does not replace Su's stronger vanishing and $\varphi(k)$-period
theorem for A281428 \cite[Theorem 4.1]{su}.

All series below are formal power series over $\QQ$, all integrals have
zero constant term, and $[t^j]$ denotes ordinary coefficient extraction.
For a positive integer $q$, let $C_q,S_q$ be the unique solutions
\begin{equation}\label{system}
 C_q'=C_q^qS_q,\qquad S_q'=C_q^{q+1},\qquad
 C_q(0)=1,\quad S_q(0)=0.
\end{equation}
Coefficients are determined successively by these equations. Differentiation
gives $C_q^2-S_q^2=1$. Thus $H_q=\log(C_q+S_q)$ satisfies
$H_q'=C_q^q$, $H_q(0)=0$, and $C_q=\cosh H_q$.
Conversely these last equations, with $S_q=\sinh H_q$, give
\eqref{system}. This identifies the defining series in A281181 when $q=3$
and A281440 when $q=7$.

\begin{samepage}
\begin{theorem}\label{main}
For every positive integer $q$ and integer $e\geq0$, the numbers
\[
 A_{q,e}(n)=(2n)![t^{2n}]C_q(t)^e\qquad(n\geq0)
\]
are integers. For all $k\geq1$ and $n\geq k^2$,
\begin{equation}\label{period}
 A_{q,e}\bigl(n+\varphi(k)^2\bigr)\equiv A_{q,e}(n)\pmod{k}.
\end{equation}
\end{theorem}
\end{samepage}

The primary conjecture is $(q,e)=(3,1)$. The other specializations
\cite{oeis180,oeis183,oeis184,oeis440} are
\begin{equation}\label{identifications}
\begin{aligned}
 \mathrm{A281183}(n)&=A_{3,2}(n),&
 \mathrm{A281184}(n)&=A_{3,3}(n),\\
 \mathrm{A281440}(n)&=A_{7,1}(n),&
 \mathrm{A281180}(n)&=A_{3,4}(n-1)\quad(n\geq1).
\end{aligned}
\end{equation}
The last identity follows from $S_3'=C_3^4$, so its sufficient bound in
the source index is $n\geq k^2+1$.

\section{A finite-jet congruence}
For a prime $p$ and $r\geq1$ put $L=p^r$ and let $I$ denote an identity
matrix. A jet of length $L$ is an element of the truncated polynomial
ring $R[x]/(x^L)$.

\begin{lemma}\label{jet}
Let $B$ be an integer $L\times L$ matrix whose reduction modulo $p$ is
the operator $P(x)\,d/dx$ on $\FF_p[x]/(x^L)$, where $P(0)=1$.
Put $U=p^{2r-1}$ and $T=p^{2r-2}(p-1)^2$. Then
\begin{equation}\label{annihilator}
 B^U(B^T-I)\equiv0\pmod{p^r}.
\end{equation}
Consequently $B^{j+T}\equiv B^j\pmod{p^r}$ for every $j\geq U$.
\end{lemma}
\begin{proof}
First work on $J=\FF_p[x]/(x^p)$ with $\delta=P\,d/dx$.
The iterated Leibniz rule shows that $\delta^p$ is a derivation: its
intermediate binomial coefficients vanish modulo $p$. Every derivation
of $J$ is determined by its value on $x$, and $P$ is a unit. Hence
$\delta^p=h\delta$ for some $h\in J$. Commuting $\delta$ with
$\delta^p$ and applying the resulting identity to $x$ gives
$\delta(h)P=0$. Therefore $h'=0$, and a polynomial of degree less
than $p$ with zero derivative is constant. Write $h=\alpha\in\FF_p$.

For $t=(p-1)^2$, this implies
\begin{equation}\label{local}
 Q(\delta):=\delta^p(\delta^t-I)=0.
\end{equation}
If $\alpha=0$ this is immediate. Otherwise repeated use of
$\delta^p=\alpha\delta$ gives
$\delta^{1+t}=\alpha^{p-1}\delta=\delta$, and multiplication by
$\delta^{p-1}$ gives \eqref{local}. This includes $p=2$.

On $\FF_p[x]/(x^L)$ consider the filtration by $(x^{jp})$, for
$0\leq j\leq p^{r-1}$. Since $\delta(x^{jp})=0$, each successive
quotient carries exactly the length-$p$ operator just considered.
Thus $Q(\delta)$ moves each filtration step into the next, and
$Q(\delta)^{p^{r-1}}=0$. Taking this power in characteristic $p$ gives
\[
 B^{p^r}\bigl(B^{(p-1)^2p^{r-1}}-I\bigr)\equiv0\pmod p.
\]
Set $u=p^r$, $v=(p-1)^2p^{r-1}$, $X=B^{u+v}$ and $Y=B^u$.
These integer matrices commute and $X\equiv Y\pmod p$.
The binomial theorem gives
$X^{p^{r-1}}\equiv Y^{p^{r-1}}\pmod{p^r}$: at each step,
raising commuting matrices congruent modulo $p^a$ to the $p$th power
raises their congruence to modulo $p^{a+1}$, for $a\geq1$.
This is \eqref{annihilator}. Multiplying it by $B^{j-U}$ proves
the final assertion.
\end{proof}

\section{Proof of Theorem \ref{main}}
On $\ZZ[c,s]$ define the integral derivation and evaluation
\[
 D(c)=c^qs,\qquad D(s)=c^{q+1},\qquad \ev(F)=F(1,0).
\]
The iterated Leibniz rule shows that
$F\mapsto\sum_{j\geq0}\ev(D^jF)t^j/j!$ preserves products.
Its values on $c,s$ consequently satisfy \eqref{system} and equal
$C_q,S_q$. Replacing $(t,s)$ by $(-t,-s)$ gives the same system,
so $C_q$ is even and $S_q$ is odd. In particular
\begin{equation}\label{integral}
 A_{q,e}(n)=\ev(D^{2n}c^e)\in\ZZ.
\end{equation}
The relation $c^2-s^2=1$ can be imposed throughout, since $D$ annihilates
$c^2-s^2-1$ and evaluation annihilates it as well.

\paragraph{Odd prime powers.}
Fix an odd prime $p$, $r\geq1$ and $L=p^r$; let $R=\ZZ/p^r\ZZ$.
On $J=R[y]/(y^L)$ there is a unique $b\in1+yJ$ with
$b^2=1+y^2$: solve successive
coefficients, dividing only by the unit $2$.
Define $\delta=b^{q+1}d/dy$; it is well defined because the
derivative of $y^L$ is zero in $R$. Its polynomial coefficient has
constant term one. When $q$ is odd this coefficient is
$(1+y^2)^{(q+1)/2}$, but the definition with $b$ works for every
positive integer $q$.
Differentiating this identity gives
$\delta(b)=yb^q$ and $\delta(y)=b^{q+1}$.
Thus the substitution $c\mapsto b$, $s\mapsto y$ intertwines $D$
with $\delta$, and evaluation at $y=0$ corresponds to $\ev$.
Apply Lemma \ref{jet} to an integer lift of the matrix of $\delta$.
For $j\geq U=p^{2r-1}$,
\[
 \ev(D^{j+T}c^e)\equiv\ev(D^jc^e)\pmod{p^r},
 \qquad T=\varphi(p^r)^2.
\]
Applying this twice with $j=2n$ shows
$A_{q,e}(n+T)\equiv A_{q,e}(n)\pmod{p^r}$ whenever $n\geq L^2$.

\paragraph{Powers of two.}
Set $z=c-1$. Applying $D$ twice to a polynomial $f(c-1)$, and using
$s^2=c^2-1$, gives the integer differential operator
\begin{equation}\label{second}
 Kf=a(z)f''+b(z)f',\qquad
 \begin{cases}
 a(z)=(1+z)^{2q}\bigl((1+z)^2-1\bigr),\\
 b(z)=(q+1)(1+z)^{2q+1}-q(1+z)^{2q-1}.
 \end{cases}
\end{equation}
Consequently $A_{q,e}(n)=(K^n(1+z)^e)(0)$.
For $L=2^r$, the ideal $(z^L)$ is preserved by $K$ modulo $L$:
the possibly surviving differentiated factors have coefficients
$L$, $2L$ or $L(L-1)$, all divisible by $L$.
Thus $K$ has a length-$L$ matrix modulo $L$.
Modulo two every polynomial has zero second derivative, so this matrix
reduces to $b(z)\,d/dz$, with $b(0)=1$.
Lemma \ref{jet}, now with $p=2$, gives
\[
 A_{q,e}\bigl(n+\varphi(2^r)^2\bigr)
 \equiv A_{q,e}(n)\pmod{2^r}\qquad(n\geq2^{2r-1}).
\]

\paragraph{An arbitrary modulus.}
For each exact prime-power factor $p^r\Vert k$, $\varphi(p^r)$ divides
$\varphi(k)$, by multiplicativity of the totient. Also
$n\geq k^2$ implies $n\geq p^{2r}$.
The preceding prime-power periods therefore divide $\varphi(k)^2$,
with a sufficient common starting index. The Chinese remainder theorem
proves \eqref{period}. For $k=1$ the congruence is automatic.
\hfill$\square$

\section{Companion formulas and exact examples}
The identifications \eqref{identifications} give all five source
conjectures. The coverage is as follows; the revision numbers identify
the cited versions.
\begin{center}\small
\begin{tabular}{lll}
\toprule
Source, revision & Column & Coverage\\
\midrule
A281181, 103 & $A_{3,1}(n)$ & Reproved; prior proof Su (2026)\\
A281183, 60 & $A_{3,2}(n)$ & Reproved; prior proof Su (2026)\\
A281184, 41 & $A_{3,3}(n)$ & Reproved; prior proof Su (2026)\\
A281180, 86 & $A_{3,4}(n-1)$ & Reproved; prior proof Su (2026)\\
A281440, 17 & $A_{7,1}(n)$ & Reproved; prior proof Su (2026)\\
A281432, 10 & $A_{5,1}(n)$ & Corollary; prior theorem Su (2026)\\
A281439, 6 & $A_{7,8}(n-1)$ & Corollary; prior theorem Su (2026)\\
\bottomrule
\end{tabular}
\end{center}
The last two entries are additional specializations
\cite{oeis432,oeis439}: the first has $H'=C^5$, and the second has
$S_7'=C_7^8$. Thus their sufficient source-index bounds are $k^2$
and $k^2+1$, respectively. Both also follow from Su's general
Theorem 3.2, with $a=q+1$ in his notation.
\begin{corollary}\label{companions}
The even-exponent companions in the following table satisfy the
square-totient congruence of Theorem \ref{main}, with sufficient
source-index bound $k^2$ for an even column and $k^2+1$ for an odd column.
The same holds for A281431. The full-index sequence A281182 satisfies
$a(j+2\varphi(k)^2)\equiv a(j)\pmod{k}$ for $j\geq2k^2$.
\end{corollary}
\begin{proof}
For every $q$, $S_q'=C_q^{q+1}$, so
$(2n-1)![t^{2n-1}]S_q=A_{q,q+1}(n-1)$ for $n\geq1$.
The defining equations in the cited entries give the table's
identifications. Apply Theorem \ref{main}.
For A281182, its even and odd indexed terms are respectively
$A_{3,1}(n)$ and $A_{3,4}(n)$ at indices $2n$ and $2n+1$;
shifting both indices by $2\varphi(k)^2$ proves the assertion.
\end{proof}
\begin{center}\small
\begin{tabular}{lll}
\toprule
Source, revision & Column & Coverage\\
\midrule
A281428, 54 & $A_{4,1}(n)$ & Not posed; proved (weaker than Su)\\
A281427, 29 & $A_{4,5}(n-1)$ & Not posed; proved\\
A281436, 10 & $A_{6,1}(n)$ & Not posed; proved\\
A281435, 19 & $A_{6,7}(n-1)$ & Not posed; proved\\
A281444, 9 & $A_{8,1}(n)$ & Not posed; proved\\
A281443, 5 & $A_{8,9}(n-1)$ & Not posed; proved\\
A281431, 6 & $A_{5,6}(n-1)$ & Corollary; prior theorem Su (2026)\\
A281182, 32 & Interlaced $C_3+S_3$ & Not posed; proved (double shift)\\
A281428, claims 1--2 & Vanishing; $\varphi(k)$-shift & Prior proof Su (2026)\\
\bottomrule
\end{tabular}
\end{center}
Here ``not posed'' refers to the square-totient (or double-shift)
congruence just stated, not to an assertion of error in the source.
The two actual conjectures in Bala's 3 March 2026 block of A281428
concern eventual vanishing for moduli coprime to six and an eventual
$\varphi(k)$-period otherwise; both have prior proofs in Su's
Theorem 4.1. They are cited, not reproved here. The source definitions
for the further companions are \cite{oeis182,oeis427,oeis428,oeis431,
oeis435,oeis436,oeis443,oeis444}.
The $q=3$ recurrence recorded by Seiichi Manyama on 14 April 2026
\cite{oeis181,oeis183,oeis184} follows directly from
\[
 D^2c^e=e(e+q)c^{e+2q}-e(e+q-1)c^{e+2q-2}.
\]
For $q=3$ this gives
$A_{3,e}(n)=e(e+3)A_{3,e+6}(n-1)-e(e+2)A_{3,e+4}(n-1)$.

One tentative formula in Bala's 22 December 2025 block of A281183 also
has a short proof, as in Su's Proposition 3.5. Put $u(y)=1+y^2$,
$\mathcal D=u(y)^2\,d/dy$ and $\mathcal Lf=u(y)(u(y)f)'$.
On $\QQ[[y]]$, multiplication by $u$ is invertible and
$\mathcal L=u^{-1}\mathcal D u$. Therefore
$\mathcal L^{2n}1(0)=\mathcal D^{2n}u(0)$.
Since $S_3'=u(S_3)^2$ and $u(S_3)=C_3^2$, this equals
$A_{3,2}(n)$, proving the formula for all $n\geq0$.

The following exact periods, previously given by Su
\cite[Proposition 3.6]{su}, illustrate the distinction between a
permitted period and the least one. The initial index specifies the
tail whose least period is stated.
\begin{center}
\begin{tabular}{lrrr}
\toprule
Sequence & Modulus & Least period & Initial index\\
\midrule
A281180 & 7 & 18 & 1\\
A281181 & 10 & 8 & 0\\
A281181 & 9 & 3 & 1\\
A281181 & 13 & 72 & 0\\
A281183 & 11 & 25 & 1\\
A281184 & 11 & 25 & 1\\
A281440 & 13 & 72 & 0\\
\bottomrule
\end{tabular}
\end{center}
These statements have finite proofs, not just matching prefixes.
For any modulus $m\geq2$, the operator \eqref{second} preserves
$(z^m)$ modulo $m$ by the same divisibility argument used above.
Let $B$ be its $m\times m$ matrix. Each scalar difference
$[z^0]B^j(B^P-I)(1+z)^e$ satisfies the monic characteristic
recurrence of $B$. Vanishing at $m$ consecutive indices proves
vanishing thereafter by Cayley--Hamilton, over $\ZZ/m\ZZ$ as well
as over fields. The supplied \texttt{verify.py} gives these zero
differences and, for each prime divisor $\ell$ of $P$, a nonzero
difference at shift $P/\ell$. They exclude every proper divisor
period and hence certify minimality.

The same calculation proves
\[
 \mathrm{A281181}(n)\equiv1\pmod{12},\qquad
 \mathrm{A281181}(n+6)\equiv7\,\mathrm{A281181}(n)\pmod{13}
 \quad(n\geq0).
\]
For these identities, respectively twelve and thirteen initial scalar
differences are zero. The second identity gives the geometric columns
in Bala's displayed modulo-$13$ array. The broader informal statement
about a two-dimensional layout is not needed for the period theorem;
Su's Proposition 3.4 supplies a precise finite-difference formulation.
Corollary \ref{companions} gives the separate double-shift statement
for the interlacing, without confusing it with one even column.

\section{Verification coverage}
The complete proof above covers all positive integer $q$, all $e\geq0$,
and the displayed bound $n\geq k^2$.
At revision \texttt{92de86888730265f84e24d28bdb67546c543a966},
the declaration \texttt{result} in each of the following Lean files
gives the source-defined series and an eventual totient-square period:
\begin{itemize}
\item \href{https://github.com/the-omega-institute/trureturing/blob/92de86888730265f84e24d28bdb67546c543a966/D5/S1/Recurrence/Periodic/BalaA281181TotientSquarePeriod.lean}{\texttt{BalaA281181TotientSquarePeriod.lean}}: A281181;
\item \href{https://github.com/the-omega-institute/trureturing/blob/92de86888730265f84e24d28bdb67546c543a966/D5/S1/Recurrence/Periodic/BalaA281183TotientSquarePeriod.lean}{\texttt{BalaA281183TotientSquarePeriod.lean}}: A281183;
\item \href{https://github.com/the-omega-institute/trureturing/blob/92de86888730265f84e24d28bdb67546c543a966/D5/S1/Recurrence/Periodic/BalaA281184TotientSquarePeriod.lean}{\texttt{BalaA281184TotientSquarePeriod.lean}}: A281184;
\item \href{https://github.com/the-omega-institute/trureturing/blob/92de86888730265f84e24d28bdb67546c543a966/D5/S1/Recurrence/Periodic/BalaA281180TotientSquarePeriod.lean}{\texttt{BalaA281180TotientSquarePeriod.lean}}: A281180, with its index offset;
\item \href{https://github.com/the-omega-institute/trureturing/blob/92de86888730265f84e24d28bdb67546c543a966/D5/S1/Recurrence/Periodic/BalaA281440TotientSquarePeriod.lean}{\texttt{BalaA281440TotientSquarePeriod.lean}}: A281440.
\end{itemize}
In the A281181 file, \texttt{BalaPolynomialSystem.even\_sequence\_period}
has the uniform scope $q\in\{3,7\}$ and every $e\geq0$.
Its statement gives an unspecified starting index; the explicit bound
and general positive-integer-$q$ scope here come from the written proof.
The source-specific least periods, the modulo-$12$ identity and the
modulo-$13$ geometric relation are certified by the finite matrices
described above. The rational-series calculation in \texttt{verify.py}
also checks coefficient normalization against \eqref{system}.

\begin{thebibliography}{99}
\bibitem{oeis181} OEIS Foundation Inc., \emph{A281181}, revision 103,
sequence by P. D. Hanna (2017); conjecture by P. Bala, 16 February 2026;
recurrence by S. Manyama, 14 April 2026.
\url{https://oeis.org/A281181}.
\bibitem{oeis180} OEIS Foundation Inc., \emph{A281180}, revision 86,
conjecture by P. Bala, 18 February 2026.
\url{https://oeis.org/A281180}.
\bibitem{oeis183} OEIS Foundation Inc., \emph{A281183}, revision 60,
conjecture by P. Bala, 16 February 2026; derivative formula,
22 December 2025. \url{https://oeis.org/A281183}.
\bibitem{oeis184} OEIS Foundation Inc., \emph{A281184}, revision 41,
conjecture by P. Bala, 16 February 2026.
\url{https://oeis.org/A281184}.
\bibitem{oeis440} OEIS Foundation Inc., \emph{A281440}, revision 17,
conjecture by P. Bala, 19 February 2026.
\url{https://oeis.org/A281440}.
\bibitem{oeis432} OEIS Foundation Inc., \emph{A281432}, revision 10,
18 January 2026. \url{https://oeis.org/A281432}.
\bibitem{oeis439} OEIS Foundation Inc., \emph{A281439}, revision 6,
1 March 2017. \url{https://oeis.org/A281439}.
\bibitem{su} L.-H. Su, \emph{Quadratic Derivations and Uniform Modular
Periods for Six Inverse-Function Sequences}, version 1, 31 October 2026,
Theorem 3.2, Corollary 3.3 and Propositions 3.4--3.6.
\url{https://doi.org/10.5281/zenodo.23835577}.
\bibitem{bala} P. Bala, \emph{The method of Graves for computing inverse
functions}, 26 January 2026, Corollary 1, pp. 2--3.
\url{https://oeis.org/A281181/a281181.pdf}.
\bibitem{oeis182} OEIS Foundation Inc., \emph{A281182}, revision 32,
3 September 2017. \url{https://oeis.org/A281182}.
\bibitem{oeis427} OEIS Foundation Inc., \emph{A281427}, revision 29,
30 December 2025. \url{https://oeis.org/A281427}.
\bibitem{oeis428} OEIS Foundation Inc., \emph{A281428}, revision 54,
conjectures by P. Bala, 3 March 2026. \url{https://oeis.org/A281428}.
\bibitem{oeis431} OEIS Foundation Inc., \emph{A281431}, revision 6,
1 March 2017.
\url{https://oeis.org/A281431}.
\bibitem{oeis435} OEIS Foundation Inc., \emph{A281435}, revision 19,
30 December 2025. \url{https://oeis.org/A281435}.
\bibitem{oeis436} OEIS Foundation Inc., \emph{A281436}, revision 10,
18 January 2026. \url{https://oeis.org/A281436}.
\bibitem{oeis443} OEIS Foundation Inc., \emph{A281443}, revision 5,
21 May 2017. \url{https://oeis.org/A281443}.
\bibitem{oeis444} OEIS Foundation Inc., \emph{A281444}, revision 9,
18 January 2026. \url{https://oeis.org/A281444}.
\end{thebibliography}
\end{document}
